\documentclass[10pt,a4paper]{amsart}
\usepackage{amsmath,amssymb,amsthm,graphicx}
% Retained minimal preamble for the Stacks Project chapter more-algebra.tex (More on Algebra).
% The chapter file pinned-original-source.tex carries only \input{preamble} on its first line, so no
% preamble text lives inside the chapter itself.  The Stacks Project uses one single shared preamble for
% every chapter; the snapshot stages it identically as sources/stacks-chow/preamble.tex and
% sources/stacks-etale-cohomology/preamble.tex.  The block below is copied mechanically, line for line,
% from sources/stacks-chow/preamble.tex lines 14-18 (the Xy-pic setup) and 154-217 (the native theorem
% declarations together with the mathematical macros).  Nothing was added, removed or redefined; the
% wrapper class and the amsmath/amssymb/amsthm/graphicx base packages are supplied by the assembler.
% Chapter: more-algebra.tex -- More on Algebra.
% For commutative diagrams we use Xy-pic
\usepackage[all]{xy}

% We use 2cell for 2-commutative diagrams.
\xyoption{2cell}

% Theorem environments.
%
\theoremstyle{plain}
\newtheorem{theorem}[subsection]{Theorem}
\newtheorem{proposition}[subsection]{Proposition}
\newtheorem{lemma}[subsection]{Lemma}

\theoremstyle{definition}
\newtheorem{definition}[subsection]{Definition}
\newtheorem{example}[subsection]{Example}
\newtheorem{exercise}[subsection]{Exercise}
\newtheorem{situation}[subsection]{Situation}

\theoremstyle{remark}
\newtheorem{remark}[subsection]{Remark}
\newtheorem{remarks}[subsection]{Remarks}

\numberwithin{equation}{subsection}

% Macros
%
\def\lim{\mathop{\mathrm{lim}}\nolimits}
\def\colim{\mathop{\mathrm{colim}}\nolimits}
\def\Spec{\mathop{\mathrm{Spec}}}
\def\Hom{\mathop{\mathrm{Hom}}\nolimits}
\def\Ext{\mathop{\mathrm{Ext}}\nolimits}
\def\SheafHom{\mathop{\mathcal{H}\!\mathit{om}}\nolimits}
\def\SheafExt{\mathop{\mathcal{E}\!\mathit{xt}}\nolimits}
\def\Sch{\mathit{Sch}}
\def\Mor{\mathop{\mathrm{Mor}}\nolimits}
\def\Ob{\mathop{\mathrm{Ob}}\nolimits}
\def\Sh{\mathop{\mathit{Sh}}\nolimits}
\def\NL{\mathop{N\!L}\nolimits}
\def\CH{\mathop{\mathrm{CH}}\nolimits}
\def\proetale{{pro\text{-}\acute{e}tale}}
\def\etale{{\acute{e}tale}}
\def\QCoh{\mathit{QCoh}}
\def\Ker{\mathop{\mathrm{Ker}}}
\def\Im{\mathop{\mathrm{Im}}}
\def\Coker{\mathop{\mathrm{Coker}}}
\def\Coim{\mathop{\mathrm{Coim}}}

% Boxtimes
%
\DeclareMathSymbol{\boxtimes}{\mathbin}{AMSa}{"02}

%
% Macros for moduli stacks/spaces
%
\def\QCohstack{\mathcal{QC}\!\mathit{oh}}
\def\Cohstack{\mathcal{C}\!\mathit{oh}}
\def\Spacesstack{\mathcal{S}\!\mathit{paces}}
\def\Quotfunctor{\mathrm{Quot}}
\def\Hilbfunctor{\mathrm{Hilb}}
\def\Curvesstack{\mathcal{C}\!\mathit{urves}}
\def\Polarizedstack{\mathcal{P}\!\mathit{olarized}}
\def\Complexesstack{\mathcal{C}\!\mathit{omplexes}}
% \Pic is the operator that assigns to X its picard group, usage \Pic(X)
% \Picardstack_{X/B} denotes the Picard stack of X over B
% \Picardfunctor_{X/B} denotes the Picard functor of X over B
\def\Pic{\mathop{\mathrm{Pic}}\nolimits}
\def\Picardstack{\mathcal{P}\!\mathit{ic}}
\def\Picardfunctor{\mathrm{Pic}}
\def\Deformationcategory{\mathcal{D}\!\mathit{ef}}
\usepackage{cleveref}
\crefname{originallemma}{original lemma}{original lemmas}
\Crefname{originallemma}{Original lemma}{Original lemmas}
\makeatletter
\newlabel{algebra-lemma-flat-going-down}{{00HS}{source}{}{sourcepointer-1}{}}
\newlabel{algebra-lemma-flat-going-down@cref}{{[originallemma][0][]00HS}{[1][0][]source}{}{}{}}
\makeatother
\crefname{originallemma}{original lemma}{original lemmas}
\Crefname{originallemma}{Original lemma}{Original lemmas}
\makeatletter
\newlabel{algebra-lemma-dimension-base-fibre-equals-total}{{00ON}{source}{}{sourcepointer-2}{}}
\newlabel{algebra-lemma-dimension-base-fibre-equals-total@cref}{{[originallemma][0][]00ON}{[1][0][]source}{}{}{}}
\makeatother
\title{Equidimensionality and catenarity descend along a flat local map of Noetherian rings}
\author{The Stacks Project Authors}
\date{stacks-project:more-algebra.tex:local-snapshot-20261012; GNU FDL 1.2 or later}
\begin{document}
\maketitle
\setcounter{subsection}{1}
\noindent\textit{Editorial note.} All delivered mathematical text is the author's own, copied from the pinned Stacks Project chapter file more-algebra.tex; no statement or proof step was rewritten, reordered or removed. The counted statement is the labeled lemma environment at source lines 32388-32397 (canonical tag 0AW4) and the counted proof is the proof environment at source lines 32399-32435, which immediately follows it and is its own proof.\par\medskip
\noindent\textit{Editorial note.} The proof refers twice to labelled items of another chapter, namely the going down property of flat maps and the dimension formula for flat local maps; both are carried as navigation stubs with their roles recorded in the reasons. The two conclusions of the lemma are derived inside the counted span.\par\medskip
\noindent\textit{Editorial note.} The lemma that follows at source lines 32437-32441 carries a different label and is not counted, and the lemma that precedes at source lines 32324-32386 is a published item of the public index with disjoint coordinates.\par\medskip
\noindent\textit{Editorial note.} The chapter file pinned-original-source.tex carries only an input of the shared preamble on its first line, so no preamble text lives inside the chapter. The retained file sources/stacks-more-algebra/retained-preamble.tex is a mechanical, line-for-line copy of the project's single shared preamble as staged in the snapshot (Xy-pic setup plus the native theorem declarations and mathematical macros). Nothing was added, removed or redefined, and the excerpt serves only as preamble, never as a span coordinate.\par\medskip
\noindent\textit{Editorial note.} The source URL is the source-level Stacks Project URL, as required when more than one card is produced from one source; the per-claim tag is carried by claim\_key and by this note. The counted statement and proof cite no bibliography entry, so this card needs no bibliography.\par\medskip
\noindent\textit{Allowed standard prerequisites (editorial).}\begin{itemize}
\item Going down for a flat local ring map: a prime of B minimal over pB lies over p. invoked at proof lines 32400-32403. Fixes the prime q over p at which the dimension computation is performed.
\item For a flat local ring map whose special fibre has dimension zero, the dimension of the source equals the dimension of the target, and more generally the dimension formula expressing dim(B) through dim(A) and the fibre. invoked at proof lines 32404-32409 and used again at proof lines 32415-32417. Produces the chain of equalities that makes dim(B over q) independent of q and gives the two dimension identities used for A.
\item A Noetherian local ring is catenary exactly when the assignment sending a prime to the dimension of the quotient is a dimension function, i.e. drops by exactly one along each immediate specialization. used at proof lines 32424-32434. Converts the formula obtained for immediate specializations into the catenarity of A.
\end{itemize}\medskip
\ifdefined\hypertarget\hypertarget{sourcepointer-1}{}\fi
\noindent\textit{Source reference.} Original lemma 00HS: sources/stacks-chow/tags.txt, source lines 654--654. The reference is the going down property of flat maps, used at proof lines 32402-32403 only to see that a prime minimal over pB lies over p; the dimension computation that follows is inside the counted span.\par\medskip
\ifdefined\hypertarget\hypertarget{sourcepointer-2}{}\fi
\noindent\textit{Source reference.} Original lemma 00ON: sources/stacks-chow/tags.txt, source lines 901--901. The reference is the dimension formula for a flat local map used at proof lines 32406-32409 and again at proof lines 32415-32417; the equalities drawn from it, and the equidimensionality and catenarity conclusions, are derived inside the counted span.\par\medskip
\begin{lemma}
\label{lemma-flat-under-catenary-equidimensional}
Let $A \to B$ be a flat local ring map of local Noetherian rings.
Assume $B$ is catenary and is $\Spec(B)$ equidimensional. Then
\begin{enumerate}
\item $\Spec(B/\mathfrak p B)$ is equidimensional for all
$\mathfrak p \subset A$ and
\item $A$ is catenary and $\Spec(A)$ is equidimensional.
\end{enumerate}
\end{lemma}

\begin{proof}
Let $\mathfrak p \subset A$ be a prime ideal. Let $\mathfrak q \subset B$
be a prime minimal over $\mathfrak pB$. Then $\mathfrak q \cap A = \mathfrak p$
by going down for $A \to B$
(Algebra, Lemma \ref{algebra-lemma-flat-going-down}).
Hence $A_\mathfrak p \to B_\mathfrak q$ is a flat local ring map
with special fibre of dimension $0$ and hence
$$
\dim(A_\mathfrak p) = \dim(B_\mathfrak q) = \dim(B) - \dim(B/\mathfrak q)
$$
(Algebra, Lemma \ref{algebra-lemma-dimension-base-fibre-equals-total}).
The second equality because $\Spec(B)$ is equidimensional and $B$ is catenary.
Thus $\dim(B/\mathfrak q)$ is independent of the choice of $\mathfrak q$
and we conclude that $\Spec(B/\mathfrak p B)$ is equidimensional of
dimension $\dim(B) - \dim(A_\mathfrak p)$. On the other hand, we
have
$\dim(B/\mathfrak p B) = \dim(A/\mathfrak p) + \dim(B/\mathfrak m_A B)$
and
$\dim(B) = \dim(A) + \dim(B/\mathfrak m_A B)$
by flatness (see lemma cited above) and we get
$$
\dim(A_\mathfrak p) = \dim(A) - \dim(A/\mathfrak p)
$$
for all $\mathfrak p$ in $A$. Applying this to all minimal primes in
$A$ we see that $A$ is equidimensional.
If $\mathfrak p \subset \mathfrak p'$ is a strict inclusion
with no primes in between, then we may apply the above to
the prime $\mathfrak p'/\mathfrak p$ in $A/\mathfrak p$
because $A/\mathfrak p \to B/\mathfrak p B$ is flat and
$\Spec(B/\mathfrak p B)$ is equidimensional, to get
$$
1 = \dim((A/\mathfrak p)_{\mathfrak p'}) =
\dim(A/\mathfrak p) - \dim(A/\mathfrak p')
$$
Thus $\mathfrak p \mapsto \dim(A/\mathfrak p)$ is a dimension
function and we conclude that $A$ is catenary.
\end{proof}

\end{document}
